\subsection{闭中心子群的情形}

在本目中，我们假设群 G 是单模的。

 考虑 G 的中心的一个闭子群 Z，并用 dz 表示 Z 上的一个哈尔测度。我们给商群 G/Z 赋以哈尔测度 \(\nu = \mu/dz\)。设 \(\chi\) 是 Z 的一个酉特征标。商群 G/Z 是单模的（INT, VII, 第 61 页, § 2, 第 7 目, 命题 11 的推论）。

引理 7. --- 我们有

\[
\mathrm{N} _ {1} (\check {f}) = \mathrm{N} _ {1} (f), \quad \langle f _ {1} | f _ {2} \rangle = \langle \check {f} _ {1} | \check {f} _ {2} \rangle .\tag{6}
\]

对所有 \(f \in \mathcal{F}_{\chi}(\mathrm{G})\) 以及一切 \(f_1\) 与 \(f_2\)（均属于 \(\mathcal{L}_{\chi}^{2}(\mathrm{G})\)）成立。

这些公式都是定义的直接推论。

 对每个 \(g \in \mathbf{G}\) 与每个 \(f \in \mathcal{F}_{\chi}(\mathbf{G})\)，函数 \(x \mapsto f(g^{-1}x)\) 与 \(x \mapsto f(xg)\) 都属于 \(\mathcal{F}_{\chi}(\mathbf{G})\)。分别把它们记作 \(\boldsymbol{\gamma}_{\mathrm{G},\chi}(g)f\) 与 \(\boldsymbol{\delta}_{\mathrm{G},\chi}(g)f\)。映射 \(\boldsymbol{\gamma}_{\mathrm{G},\chi}\) 与 \(\boldsymbol{\delta}_{\mathrm{G},\chi}\) 是 \(\mathbf{G}\) 在 \(\mathcal{F}_{\chi}(\mathbf{G})\) 中的线性表示。对每个 \(z \in \mathbf{Z}\)，我们有

\[
\boldsymbol {\gamma} _ {\mathrm{G}, \chi} (g z) = \overline {{\chi (z)}} \boldsymbol {\gamma} _ {\mathrm{G}, \chi} (g), \quad \boldsymbol {\delta} _ {\mathrm{G}, \chi} (g z) = \chi (z) \boldsymbol {\delta} _ {\mathrm{G}, \chi} (g).\tag{7}
\]

设 \(f\in \mathcal{F}_{\chi}(\mathrm{G})\) 且 \(g\in \mathrm{G}\)；我们有

\[
| \boldsymbol {\gamma} _ {\mathrm{G}, \chi} (g) f | = \boldsymbol {\gamma} _ {\mathrm{G} / \mathrm{Z}} (g \mathrm{Z}) | f |, \quad | \boldsymbol {\delta} _ {\mathrm{G}, \chi} (g) f | = \boldsymbol {\delta} _ {\mathrm{G} / \mathrm{Z}} (g \mathrm{Z}) | f |,\tag{8}
\]

其中出现在这些等式中的所有函数都等同于 G/Z 上的函数。

 子空间 \(\mathcal{K}_{\chi}(\mathrm{G})\)（\(\mathcal{F}_{\chi}(\mathrm{G})\) 的）在 \(\gamma_{\mathrm{G},\chi}\) 与 \(\delta_{\mathrm{G},\chi}\) 下稳定。设 \(p\) 是一个实数 \(\geqslant 1\)。公式 (8) 蕴含：表示 \(\gamma_{\mathrm{G},\chi}\) 与 \(\delta_{\mathrm{G},\chi}\) 在限制到 \(\mathcal{K}_{\chi}(\mathrm{G})\) 上后可延拓为 G 到 \(\mathcal{L}_{\chi}^{p}(\mathrm{G})\) 中的连续等距线性表示，分别记作 \(\gamma_{\mathrm{G},\chi}^{(p)}\) 与 \(\delta_{\mathrm{G},\chi}^{(p)}\)。通过取商，这些表示也在 \(\mathrm{L}_{\chi}^{p}(\mathrm{G})\) 中定义了 G 的等距表示，仍以同样的记号记之。

 表示 \(\gamma_{\mathrm{G},\chi}^{(2)}\) 与 \(\delta_{\mathrm{G},\chi}^{(2)}\)（\(\mathrm{L}_{\chi}^{2}(\mathrm{G})\) 中的）是酉表示，将分别简记为 \(\gamma_{\mathrm{G},\chi}\) 与 \(\delta_{\mathrm{G},\chi}\)——当不致与 \(\mathcal{F}_{\chi}(\mathrm{G})\) 中的表示混淆时。我们也用 \(\varrho_{\mathrm{G},\chi}\) 表示 \(\mathrm{G} \times \mathrm{G}\) 在 \(\mathcal{L}_{\chi}^{2}(\mathrm{G})\) 或 \(\mathrm{L}_{\chi}^{2}(\mathrm{G})\) 中的由下式定义的连续表示

\[
\pmb {\varrho} _ {\mathrm{G}, \chi} (g, h) = \pmb {\gamma} _ {\mathrm{G}, \chi} (g) \circ \pmb {\delta} _ {\mathrm{G}, \chi} (h) = \pmb {\delta} _ {\mathrm{G}, \chi} (h) \circ \pmb {\gamma} _ {\mathrm{G}, \chi} (g)
\]

对所有 \((g,h)\in \mathrm{G}\times \mathrm{G}\) 成立（参见 V, 第 377 页, 引理 1）。

\(\gamma_{\mathrm{G},\chi}\) 在 \(\mathrm{L}_{\chi}^{2}(\mathrm{G})\) 中的表示不是别的，正是诱导表示 \(\mathrm{Ind}_{\mathrm{Z}}^{\mathrm{G}}(\overline{\chi})\)。

当 \(Z = \{e\}\) 时，由于 G 单模，下面的引理可由 INT, VIII, 第 166 页, § 4, 第 5 目, 命题 12 得出。

引理 8. --- 设 \(f_1 \in \mathcal{K}_\chi(\mathrm{G})\) 且 \(f_2 \in \mathcal{L}_\chi^2(\mathrm{G})\)。

\begin{enumerate}
\def\labelenumi{\alph{enumi})}
\item
  函数 \(f\)——它定义在 \(\mathrm{G}\) 上，由 \(f(g) = \langle f_1 | \boldsymbol{\gamma}_{\mathrm{G},\chi}(g)f_2 \rangle\)（对所有 \(g \in \mathrm{G}\)）给出——属于 \(\mathcal{L}_{\chi}^{2}(\mathrm{G})\)，且满足 \(\mathrm{N}_2(f) \leqslant \mathrm{N}_1(f_1)\mathrm{N}_2(f_2)\)；
\item
  函数 \(f\)——它定义在 \(G\) 上，由 \(f(g) = \langle f_1 | \delta_{G,\chi}(g)f_2 \rangle\)（对所有 \(g \in G\)）给出——属于 \(\mathcal{L}_{\chi}^{2}(G)\)，且满足 \(N_2(f) \leqslant N_1(f_1)N_2(f_2)\)。
\end{enumerate}

我们来证明 \(a\) )，断言 \(b\) ) 的证明是类似的。函数 \(f\) 连续，因此 \(\mu\) -可测。对每个 \(z \in \mathbf{Z}\) 与每个 \(g \in \mathbf{G}\)，我们有

\[
f (g z) = \langle f _ {1} | \gamma_ {\mathrm{G}, \chi} (g z) f _ {2} \rangle = \overline {{\chi (z)}} \langle f _ {1} | \gamma_ {\mathrm{G}, \chi} (g) f _ {2} \rangle
\]

（公式 (7), 第 417 页），因此 \(f \in \mathcal{F}_{\overline{\chi}}(\mathrm{G})\)。V 第 413 页的命题 5 蕴含现在只需证明 \(\mathrm{N}_2(f) \leqslant \mathrm{N}_1(f_1)\mathrm{N}_2(f_2)\)。

先设 \(f_2\) 属于 \(\mathcal{K}_{\chi}(\mathrm{G})\)。对每个 \(g \in \mathrm{G}\)，按定义有

\[
f (g) = \int_ {\mathrm{G/Z}} \overline {{f _ {1}}} \gamma_ {\mathrm{G}, \chi} (g) f _ {2} d \nu ,
\]

其中函数 \(\overline{f_1} \gamma_{\mathrm{G},\chi}(g)f_2\) 等同于 G/Z 上的一个函数。

我们定义一个函数 \(f_3\)（在 \(G\) 上）：令 \(f_3(g) = 0\)（若 \(f_1(g) = 0\)），否则令 \(f_3(g) = f_1(g)|f_1(g)|^{-1/2}\)。函数 \(f_3\) 属于 \(\mathcal{F}_{\chi}(G)\) 且满足 \(f_1 = |f_1|^{1/2}f_3\)；它是 \(\mu\) -可测且模 \(Z\) 在一个紧集外为零的，因为 \(f_1\) 如此。由于 \(|f_1|^{1/2} \in \mathcal{K}_1(G)\)，由此得出

\[
\overline {{f _ {1}}} \gamma_ {\mathrm{G}, \chi} (g) f _ {2} = | f _ {1} | ^ {1 / 2} \overline {{f _ {3}}} \gamma_ {\mathrm{G}, \chi} (g) f _ {2},
\]

\[
| \overline {{f _ {3}}} \gamma_ {\mathrm{G}, \chi} (g) f _ {2} | = | f _ {3} | | \gamma_ {\mathrm{G}, \chi} (g) f _ {2} |
\]

对所有 \(g \in \mathbf{G}\) 成立。

 设 \(g \in G\)。由柯西--施瓦茨不等式，我们有

\[
\begin{array}{l} | f (g) | ^ {2} \leqslant \Big (\int_ {\mathrm{G/Z}} | f _ {1} | ^ {1 / 2} | f _ {3} | | \boldsymbol {\gamma} _ {\mathrm{G}, \chi} (g) f _ {2} | d \nu \Big) ^ {2} \\ \leqslant \Big (\int_ {\mathrm{G/Z}} | f _ {1} | d \nu \Big) \Big (\int_ {\mathrm{G/Z}} | f _ {3} | ^ {2} | \boldsymbol {\gamma} _ {\mathrm{G}, \chi} (g) f _ {2} | ^ {2} d \nu \Big). \end{array}
\]

由于有 \(|f_{3}|^{2}=|f_{1}|\)（在 \(\mathcal{K}(\mathrm{G}/\mathrm{Z})\) 中），我们通过在 G/Z 上积分得到

\[
\int_ {\mathrm{G} / \mathrm{Z}} | f (g) | ^ {2} d \nu (g) \leqslant \mathrm{N} _ {1} (f _ {1}) \int_ {\mathrm{G} / \mathrm{Z}} \left(\int_ {\mathrm{G} / \mathrm{Z}} | f _ {1} (x) | | f _ {2} (g ^ {- 1} x) | ^ {2} d \nu (x)\right) d \nu (g).
\]

由函数

\[
(g, x) \mapsto | f _ {1} (x) | | f _ {2} (g ^ {- 1} x) | ^ {2}
\]

通过取商所得到的 G/Z × G/Z 上的函数是 \((\nu \otimes \nu)\) 可测且具紧支撑的，从而是 \((\nu \otimes \nu)\) -缓增的（INT, V, 第 4 页, § 5, 第 2 目, 定义 2）。由 INT, V, 第 93 页, § 8, 第 3 目, 命题 7 得

\[
\begin{array}{l} \int_ {\mathrm{G} / \mathrm{Z}} \Bigl (\int_ {\mathrm{G} / \mathrm{Z}} | f _ {1} (x) |   | f _ {2} (g ^ {- 1} x) | ^ {2} d \nu (x) \Bigr) d \nu (g) \\ = \int_ {\mathrm{G} / \mathrm{Z}} | f _ {1} (x) | \Bigl (\int_ {\mathrm{G} / \mathrm{Z}} | f _ {2} (g ^ {- 1} x) | ^ {2} d \nu (g) \Bigr) d \nu (x) = \mathrm{N} _ {1} (f _ {1}) \mathrm{N} _ {2} (f _ {2}) ^ {2}. \end{array}
\]

于是我们有 \(\mathrm{N}_{2}(f)^{2}\leqslant\mathrm{N}_{1}(f_{1})^{2}\mathrm{N}_{2}(f_{2})^{2}\)，这就在此情形下建立了所要的性质。

 我们来考虑一般情形。用 \(u\) 表示 \(\mathcal{K}_{\chi}(\mathrm{G})\) 到 \(\mathcal{L}_{\overline{\chi}}^{2}(\mathrm{G})\) 中的把 \(f_{2}\) 映为 \(f\) 的线性映射。设 \(f_{2} \in \mathcal{L}_{\chi}^{2}(\mathrm{G})\)，并设 \((f_{2,n})_{n \in \mathbf{N}}\) 是 \(\mathcal{K}_{\chi}(\mathrm{G})\) 中的收敛于 \(f_{2}\)（在 \(\mathcal{L}_{\chi}^{2}(\mathrm{G})\) 中）的序列。令 \(f_{n} = u(f_{2,n})\)；序列 \((f_{n})_{n \in \mathbf{N}}\) 在 \(\mathcal{L}_{\overline{\chi}}^{2}(\mathrm{G})\) 中是柯西序列，因为前一种情形蕴含 \(\mathrm{N}_{2}(f_{n} - f_{m}) \leqslant \mathrm{N}_{1}(f_{1})\mathrm{N}_{2}(f_{2,n} - f_{2,m})\) 对一切 \(n\) 与 \(m\)（\(\mathbf{N}\) 中的）成立。设 \(f \in \mathcal{L}_{\overline{\chi}}^{2}(\mathrm{G})\) 使得 \((f_{n})\) 收敛于 \(f\)（V, 第 409 页, 命题 3, \(c\) )）。由于 \(\mathrm{N}_{2}(f_{n}) \leqslant \mathrm{N}_{1}(f_{1})\mathrm{N}_{2}(f_{2,n})\) 对所有 \(n \in \mathbf{N}\) 成立，由此得出 \(\mathrm{N}_{2}(f) \leqslant \mathrm{N}_{1}(f_{1})\mathrm{N}_{2}(f_{2})\)。

存在一个子序列 \((f_{n_{k}})_{k\in\mathbf{N}}\)，使得 \(f_{n_{k}}(g)\) 收敛于 \(f(g)\)，对所有位于 G 中模 Z 零测子集之外的 \(g\in G\) 成立（出处同上, b)）。但另一方面，对每个 \(g\in G\)，我们有

\[
f _ {n _ {k}} (g) = \left\langle f _ {1} \mid \gamma_ {\mathrm{G}, \chi} (g) f _ {2, n _ {k}} \right\rangle\rightarrow \left\langle f _ {1} \mid \gamma_ {\mathrm{G}, \chi} (g) f _ {2} \right\rangle .
\]

于是我们有 \(f(g) = \langle f_1 | \gamma_{\mathrm{G},\chi}(g)f_2 \rangle\)，对所有位于 G 中模 Z 零测子集之外的 \(g\) 成立。由于 \(f \in \mathcal{L}_{\overline{\chi}}^2(\mathrm{G})\) 且 \(\mathrm{N}_2(f) \leqslant \mathrm{N}_1(f_1)\mathrm{N}_2(f_2)\)，引理得证。

